% SOURCE_AUTHORITY: sga5_fr_workpass.tex, lines 11729--11769 (LF-normalized unit).
% SOURCE_SHA256: 791F4EFFC5E02832D5D77ED03518C8156D6F07E4C8238B03545DB93D883FBB28
\subsection*{3. Caráter funtorial}

Sejam $\cC$ e $\cC'$ duas categorias trianguladas, e seja $T:\cC\to\cC'$ um funtor exato de $\cC$ em $\cC'$, isto é, um funtor aditivo, «graduado» e que transforma os triângulos distinguidos em triângulos distinguidos. Denotemos, como antes, por
\[
A,A':\Ab\longrightarrow\Ab
\]
os funtores associados às categorias $\cC$ e $\cC'$. Resulta um morfismo funtorial $A'\to A$ de «composição com $T$» e, portanto, um homomorfismo
\[
K(T):K(\cC)\longrightarrow K(\cC')
\]
único tal que o diagrama
\[
\begin{tikzcd}[column sep=large,row sep=large]
\Ob\cC \arrow[r,"T"] \arrow[d,"\operatorname{cl}_{\cC}"']&
\Ob\cC' \arrow[d,"\operatorname{cl}_{\cC'}"]\\
K(\cC) \arrow[r,"K(T)"']& K(\cC')
\end{tikzcd}
\]
seja comutativo. Resulta que, se os funtores $T,T_1:\cC\to\cC'$ são isomorfos, então $K(T)=K(T_1)$. Em particular, se o funtor $T$ é uma equivalência de categorias, então $K(T)$ é um isomorfismo.

Se $\cC''$ é uma terceira categoria triangulada e
\[
F:\cC\times\cC'\longrightarrow \cC''
\]
um bifuntor exato em cada variável, então é claro que resulta uma aplicação bilinear
\[
K(F):K(\cC)\times K(\cC')\longrightarrow K(\cC''),
\]
tal que se verifica a relação seguinte:
\[
K(F)(\operatorname{cl}_{\cC}(X),\operatorname{cl}_{\cC'}(X'))
=
\operatorname{cl}_{\cC''}(F(X,X')).
\]

\begin{statement}{Proposição 3.1}
Seja $\cA$ uma categoria triangulada, $\cB$ uma subcategoria triangulada espessa de $\cA$, $\cA/\cB$ a categoria quociente de $\cA$ por $\cB$, $Q:\cA\to\cA/\cB$ o funtor canônico de passagem ao quociente, e $i:\cB\to\cA$ o funtor de inclusão. Então a sequência seguinte é exata:
\[
K(\cB)\xrightarrow{K(i)}K(\cA)\xrightarrow{K(Q)}K(\cA/\cB)\longrightarrow 0.
\]
\end{statement}
